\section{Recommendations}
Based on our review, we provide recommendations to strengthen the construct validity of LLM benchmarks. We prioritise broadly applicable recommendations, noting that not all apply to every benchmark. Each includes a specific checklist of items to address when creating a new benchmark.

\subsection{Define the phenomenon}
%Standard recommendation structure:

\begin{checklistbox}{Define the phenomenon}
    \begin{itemize}[leftmargin=5pt, label=\scriptsize$\square$]
        %\item The phenomenon being measured is clearly defined.
        \item Provide a precise and operational definition for the phenomenon being measured
        %\item The benchmark captures the entire targeted phenomenon, not a single aspect of it.
        \item Specify the scope of the phenomenon being covered and acknowledge any excluded aspects
        %\item When the benchmark aggregates sub-aspects of the phenomenon, they are described and measured separately.
        \item Identify if the phenomenon has sub-components and ensure they are measured separately
    \end{itemize}
\end{checklistbox}

%What is the desirable feature of benchmarks
The target phenomenon should be clearly defined before operationalising it in a benchmark. 78.2\% of reviewed benchmarks provide definitions, helping users to understand what is being measured and how. \textcite{phuongEvaluatingFrontierModels2024a} shows an example of how a clear definition can guide task design and clarify the operationalisation of an abstract concept.

% What can be done to improve?
When multiple definitions exist for a single term, stating one helps clarify the intention of the benchmark and how the results should be interpreted. If there is no consensus for defining a phenomenon, as we observe in 47.8\% of benchmarks, providing a negative or apophatic definition can help set boundaries (e.g., repeating memorised answers is not reasoning)~\cite{beanLINGOLYBenchmarkOlympiadlevel2024a}. If the target phenomenon has sub-components, as in 61.2\% of benchmarks, benchmarking each element separately increases clarity and improves interpretation. Although, benchmarks which combine several different measures of the same concept can be useful in no single measure is adequate.

\subsection{Measure the phenomenon and only the phenomenon}

\begin{checklistbox}{Measure only the phenomenon}
    \begin{itemize}[leftmargin=5pt, label=\scriptsize$\square$]
%\item Auxiliary tasks are either tested independently or their impact on the primary task is accounted for.
%\item Isolate or control for the impact of auxiliary tasks not central to the primary task
\item Control for unrelated tasks that may affect the results
%\item If the benchmark has a strict output format, model performance is compared with and without format constraints.
\item Assess the impact of format constraints on model performance
%\item If the benchmark uses automated techniques to parse answers, they are tested for bias and accuracy across different models.
\item Validate any automated output parsing techniques for accuracy, consistency and bias
\end{itemize}
\end{checklistbox}


%What is the desirable feature of benchmarks
Completing a benchmark task involves a combination of task-specific and more general abilities, such as instruction-following. Additional, unmeasured, subtasks can confound the measurement of the target phenomenon. For example, 21.1\% of benchmarks require specific output formats that can themselves be challenging for models~\cite{tamLetMeSpeak2024a}. Others may involve complex instructions that disproportionately reduce performance in weaker models~\cite{beanLINGOLYBenchmarkOlympiadlevel2024a}. In tasks such as commonsense reasoning, it can be difficult to separate reasoning ability from model's existing knowledge \cite{hoWikiWhyAnsweringExplaining2023}.

% What can be done to improve?
Several strategies can mitigate these confounding effects. Baselines can be established for performance on the relevant subtasks alone. If a benchmark requires world knowledge but does not intend to measure it, models should first be tested on this world knowledge directly and scores adjusted to avoid penalising failures arising from lack of knowledge. If a benchmark uses strict formats or complex instructions, test those skills independently and allow retries to distinguish formatting proficiency from task performance. If LLMs are used to parse original model responses, the extractor LLM should be validated to avoid introducing new biases or performance artifacts. Though less applicable on individual benchmarks, factor analysis techniques are also being explored to extract latent capability dimensions with less interference from auxiliary tasks \cite{ruan2024observational, ilic2024evidence, burnell2023revealing}.


\subsection{Construct a representative dataset for the task}
%Standard recommendation structure:
\begin{checklistbox}{Construct a representative dataset for the task}
    \begin{itemize}[leftmargin=5pt, label=\scriptsize$\square$]
 %\item Task items are chosen to be representative of the overall task space.
 \item Employ sampling strategies to ensure task items are representative of the overall task space
 %\item Checks are performed to ensure that the task items are high quality and related to the phenomenon (especially if the benchmark is large).
 \item Verify the quality and relevance of all task items especially for large or automatically generated datasets
 %\item The selection of task items has been tailored for testing LLMs (e.g. including both easy and difficult tasks by human standards, including small syntactic perturbations of task items).
 \item Include task items that test known LLM sensitivities (e.g. input permutations or variations)
\end{itemize}
\end{checklistbox}

%What is the desirable feature of benchmarks
Benchmarks use finite sets of task items as proxies for complex phenomena. Each item can be seen as drawn from a larger possible set, so sampling should be representative of the task space. However, 27.0\% of reviewed benchmarks used convenience sampling, relying on the validity of the existing sample. For example, if a benchmark reuses questions from a calculator-free exam such as AIME~\cite{whiteLiveBenchChallengingContaminationLimited2025}, numbers in each problem will have been chosen to facilitate basic arithmetic. Testing only on these problems would not predict performance on larger numbers, where LLMs struggle.

% What can be done to improve?
We recommend that authors adopt more robust sampling techniques, such as random or stratified sampling (17.1\% of reviewed benchmarks use at least one of these). With better sampling methods, smaller well-designed datasets can provide higher construct validity than larger datasets at less computational cost~\cite{maiapoloTinyBenchmarksEvaluatingLLMs2024}. The risk of having non-representative sampling of benchmark tasks should also be taken into account when generating synthetic examples to increase the size of benchmarks, as occurs in 47.5\% of the benchmarks we reviewed. Task items can also be explicitly designed to test for common weaknesses in LLMs. For example, human examinations are unlikely to have the same question repeated in several different phrasings, but LLMs are known to be sensitive to minor variations in prompts \cite{zhuoProSAAssessingUnderstanding2024, mizrahiStateWhatArt2024} and variations could improve the robustness of the results.

\subsection{Acknowledge limitations of reusing datasets}
%Standard recommendation structure:
\begin{checklistbox}{Acknowledge limitations of reusing datasets}
    \begin{itemize}[leftmargin=5pt, label=\scriptsize$\square$]
    %\item The benchmark documents whether it adapts a previous dataset or benchmark.
    \item Document whether the benchmark adapts a previous dataset or benchmark
    %\item If so, the authors provide a clear analysis of the previous work describing strengths and limitations of the benchmark.
    \item If so, analyse and report the relevant strengths and limitations of the adapted prior work
    %\item If so, results on the original benchmark are included and compared to.
    \item If so, report and compare performance on the new benchmark against the original
    %\item Any differences from the original dataset are justified and explained in the context of construct validity.
    %\item Articulate the rationale when modifying a reused dataset, identifying improvements to construct validity
    \item Explain modifications to reused datasets and how they improve construct validity
\end{itemize}
\end{checklistbox}

%What is the desirable feature of benchmarks
38.2\% of reviewed benchmarks reuse data from previous benchmarks or human exams. Reusing existing datasets makes it difficult for authors to control the construct validity of their benchmark and limits options for design choices such as task and metric. Reuse of existing materials also increases the chances of benchmarking tasks appearing in pre-training data (see \Cref{rec:contamination}), compromising results.

% What can be done to improve?

Newly constructed datasets should be preferred to reused datasets.
When datasets are reused, such as when a benchmark improves upon an older version (7.7\% of cases), authors must investigate which changes have been introduced and what the new benchmark preserves.
Differences between the original and new datasets should be clearly documented and justified. Reporting differences in results between the new and original benchmarks can help to demonstrate the impact of changes, including whether the new benchmark has improved the construct validity.

\subsection{Prepare for contamination}
\label{rec:contamination}
\begin{checklistbox}{Prepare for contamination}
    \begin{itemize}[leftmargin=5pt, label=\scriptsize$\square$]
    %\item The benchmark includes a method for identifying contamination.
    \item Implement tests to detect data contamination and apply them to the benchmark
    %\item Held-out task items are available (e.g. via a private test set, or generating new questions).
    \item Maintain a held-out set of task items to facilitate ongoing, uncontaminated evaluation
    %\item The authors conduct an analysis of data exposure prior to the benchmark creation.
    \item Investigate the potential pre-exposure of benchmark source materials or similar data in common LLM training corpora
\end{itemize}
\end{checklistbox}

For many phenomena, the process through which an LLM reaches the answer is equally as important as whether the correct answer was reached. In these cases, the validity of the results can be undermined both by direct contamination of benchmark items and by memorisation of partial answers or closely-related information. Benchmark contamination is likely to occur even with model developers acting in good faith \cite{zhouDontMakeYour2023}. 
When a benchmark is widely used, the progressive effort to improve on that benchmark means that newly developed techniques will be specifically suited to solving that task. Over time, this selection of methods can lead to overfitting, similar to the repeated use of a validation set \cite{zhangCarefulExaminationLarge2024a}, effectively contaminating the benchmark.

We recommend vetting test items for dataset contamination when the benchmark is created, especially when the dataset is already public, or when an LLM is used to generate task examples \cite{mainiLLMDatasetInference2024}. Including these contamination checks within the benchmark itself can provide ongoing verification of the validity. Dynamic benchmarks have also been proposed as a solution to this issue \cite{kielaDynabenchRethinkingBenchmarking2021d}, and procedurally generated tasks, in particular, can be used to keep benchmarks up-to-date \cite{khoujaLINGOLYTOODisentanglingMemorisation2025}.


\subsection{Use statistical methods to compare models}
\begin{checklistbox}{Use statistical methods to compare models}
    \begin{itemize}[leftmargin=5pt, label=\scriptsize$\square$]
     %\item The sample size of the benchmark is reported and sufficiently powered.
     \item Report the benchmark's sample size and justify its statistical power
     %\item Uncertainty estimates are provided for the main scores, and are narrow enough to meaningfully compare relevant models.
     \item Report uncertainty estimates for all primary scores to enable robust model comparisons
     %\item If human raters are used, the recruitment accounts for demographic biases which may be relevant to the preferences they report.
     \item If using human raters, describe their demographics and mitigate potential demographic biases in rater recruitment and instructions
    %\item If the benchmark uses subjective measures of performance, the distribution of labels is reported and accounted for in the scoring.
    \item Use metrics that capture the inherent variability of any subjective labels, without relying on single-point aggregation or exact matching

\end{itemize}
\end{checklistbox}

To support valid interpretation and comparisons across models, prior work has highlighted the importance of using statistical techniques in the analysis of benchmark results~\cite{millerAddingErrorBars2024,mcintoshInadequaciesLargeLanguage2024}. At present, only 16.0\% of reviewed benchmarks conducted any statistical testing. Increasing the use of robust statistical methods for LLM benchmarking is critical.

In addition, scoring methods based on human or LLM ratings provide subjective metrics that may vary across samples. Since there is real variation in human preferences, the aggregation and reporting should consider the meaning of the distribution of ratings. In particular, benchmark creators should consider the representativeness of the raters, and whether there are meaningful differences between groups \cite{kirkPRISMAlignmentDataset2024}. For example, bias benchmarks operate with concepts of harm and bias which are culturally and socially contingent \cite{sahooIndiBiasBenchmarkDataset2024}. By considering the distribution of ratings, overall results will better incorporate potential real-world uses of LLMs.

\subsection{Conduct an error analysis}
\begin{checklistbox}{Conduct an error analysis}
    \begin{itemize}[leftmargin=5pt, label=\scriptsize$\square$]
    %\item The authors include an analysis of failure cases.
    \item Conduct a qualitative and quantitative analysis of common failure modes
    %\item There are no patterns of failures which relate to non-targeted phenomena.
    \item Investigate whether failure modes correlate with non-targeted phenomena (confounders) rather than the intended construct
    %\item If so, potential biases in the scoring are discussed.
    \item If so, identify and discuss any potential scoring biases revealed in the error analysis
    %\item Experiments or recommendations are made to improve model scores on the benchmark.
    %\item Conduct experiments or propose new directions to improve model scores on the benchmark
\end{itemize}
\end{checklistbox}

After a benchmark is created, an error analysis can reveal the types of errors models make. If the benchmark has high construct validity, these errors will indicate useful research directions for the target phenomenon. Therefore, error analysis can provide an indication of the construct validity of the benchmark based on the avenues for improvement which are indicated. If the failure cases correspond to failures to demonstrate the target phenomenon, the validity is high. If not, this may be a reason to modify the benchmark to be more precise. As an example, \citeauthor{phuongEvaluatingFrontierModels2024a} experiment with repeated trials and find that the highest scores come from low probability generations, allowing them to identify that the tasks are possible for the models, but not likely to be solved.

\subsection{Justify construct validity}
\vspace{-5pt}
\begin{checklistbox}{Justify construct validity}
    \begin{itemize}[leftmargin=5pt, label=\scriptsize$\square$]
    %\item A relevant use case is provided to justify the direction of research.
    \item Justify the relevance of the benchmark for the phenomenon with real-world applications
    %\item The authors explain why the task and metric were chosen to measure the target phenomenon.
    \item Provide a clear rationale for the choice of tasks and metrics, connected to the operational definition of the phenomenon
    %\item The benchmark is compared to other benchmarks of similar phenomena, with a discussion of similarities and differences.
    \item Compare similarities and differences between the benchmark and existing evaluations of similar phenomena
    \item Discuss the limitations and design trade-offs of the benchmark concerning construct validity
\end{itemize}
\end{checklistbox}

Improving scores on a benchmark requires attention and resources, so it is helpful for the authors to explain why the benchmark is relevant. However, only about half (53.4\%) of reviewed benchmarks justify why they are a valid measure of an important phenomenon. To establish high construct validity, we recommend authors articulate the rationale behind the chain of decisions from defining a phenomenon, to operationalising it via a task, to selecting specific task items to test, to the code implementation of the task, up to making validity claims. 

Authors should discuss key design trade-offs. For example, multiple-choice formats are easy to score but can be gamed and rarely reflect real-world use cases~\cite{mccoyRightWrongReasons2019b}. Free-text responses are more realistic but harder and costlier to evaluate~\cite{bidermanLessonsTrenchesReproducible2024a}. With many benchmarks aiming at similar phenomena (e.g. `reasoning'), clarity about how a benchmark aligns or diverges from others is critical. Convergent and discriminant validity help clarify what each benchmark actually tests~\cite{alaaMedicalLargeLanguage2025a}. These choices should be addressed directly in the limitations to enable more reliable and interpretable progress.
